\section{Thin rank-three geometries and involution triples}
Throughout, the group preserves types and acts transitively on flags of
each type. Incidence is symmetric and reflexive, with no incident distinct
objects of one type. Every flag extends to a chamber. A geometry is thin
if every flag missing one type has exactly two extensions. Residual
connectedness requires the incidence graph of every residue of rank at
least two to be connected. Residual weak primitivity means that in every
nonempty residue at least one of its actions on a type is primitive.

\begin{lemma}\label{thin-triples}
Let a nonabelian finite simple group $G$ act as above on a thin, residually
connected geometry of rank three. For a chamber $C=\{p_0,p_1,p_2\}$,
write $H_i=G_{p_i}$ and $B=\bigcap_i H_i$. There are distinct involutions
$t_0,t_1,t_2$ such that
\[
 B=1,\qquad G=\langle t_0,t_1,t_2\rangle,\qquad
 H_i=\langle t_j:j\ne i\rangle,
 \qquad H_i\cap H_j=\langle t_k\rangle\quad(\{i,j,k\}=\{0,1,2\}).
\]
\end{lemma}
\begin{proof}
Put $R_i=\bigcap_{j\ne i}H_j$. Flag transitivity on the two extensions
of $C\setminus\{p_i\}$ gives $[R_i:B]=2$, so $B$ is normal in each $R_i$.

The chamber graph is connected. To see this, a rank-two residue has a
connected incidence graph, whose edges are its chambers; successive edges
sharing a vertex give a gallery. In rank three take a path of incident
objects between a vertex of one chamber and a vertex of another. Extend
each edge to a chamber. Consecutive chosen chambers contain a common
object, and the rank-two residue of that object supplies a gallery between
them. The end chambers connect in the same way. Thus all chambers are
connected by changes of one object.

The chambers adjacent to $C$ across its $i$-panel are images of $C$ under
$R_i$. Transporting along a gallery proves
$G=\langle R_0,R_1,R_2\rangle$: the generated subgroup reaches every
chamber, and it contains their common base stabilizer $B$. Therefore
$B\lhd G$. Since $B<R_i$, simplicity gives $B=1$. Each $R_i$ now has a
unique nonidentity element $t_i$, an involution. Within the residue of
$p_i$, the same gallery argument gives
$H_i=\langle R_j:j\ne i\rangle$. Distinctness follows because an element
in two $R_i$ lies in $B$. The stated intersections follow from the
stabilizer definitions.
\end{proof}

\begin{lemma}\label{thin-recognize}
Conversely take three distinct involutions generating $G$, and put
$H_i=\langle t_j:j\ne i\rangle$. Suppose the three intersections in
Lemma~\ref{thin-triples} hold. Their coset incidence system is a flag-transitive
geometry if and only if, for each ordering $i,j,k$ of the types,
\begin{equation}\label{products}
        (H_iH_j)\cap(H_iH_k)=H_i(H_j\cap H_k).
\end{equation}
It then is thin and residually connected. It is residually weakly primitive
if and only if at least one $H_i$ is maximal in $G$ and all three numbers
$|H_i|/2$ are prime.
\end{lemma}
\begin{proof}
Any incident pair of cosets has a common element, so $G$ is transitive
on incident pairs. Normalize the cosets of types $j,k$ to $H_j,H_k$.
The cosets $H_i g$ incident with both correspond exactly to the elements
of $(H_iH_j)\cap(H_iH_k)$. They extend to chambers in the base chamber's
orbit precisely when they meet $H_j\cap H_k$. This gives
\eqref{products}. Every flag then extends to such a chamber. The right
side is contained in the left side; since the triple intersection is
trivial, its size is $2|H_i|$. Thus equality is checked without constructing
an entire incidence graph.

The base chamber stabilizer is trivial and its panel stabilizers have
order two, so the geometry is thin. Each $H_i$ is generated by the two
panel stabilizers within it; the whole group is generated by all three.
Their coset incidence graphs are connected, proving residual connectedness.
Each rank-two residue is the polygon associated to a dihedral group of
order $2m_i$, where $m_i=|H_i|/2$. Its vertex stabilizers are reflections.
A reflection subgroup is maximal in a dihedral group of order $2m_i$
exactly when $m_i$ is prime: intermediate subgroups correspond to proper
divisors of the rotation order. Rank-one residues have degree two and
are primitive. The remaining condition is primitivity on some type of
the whole geometry, equivalent to maximality of its stabilizer $H_i$.
\end{proof}

